\chapter{Kết luận và hướng phát triển}

Số đếm syscall theo cửa sổ 0,5 s đủ để phát hiện hành vi sau xâm nhập trong container với p95 dưới 2 s trên cụm thật. Tuy vậy, chúng chưa đủ để phân biệt ác ý với một process vô hại mới.

\section{Trả lời các câu hỏi nghiên cứu}

\begin{longtable}[]{@{}>{\raggedright\arraybackslash}p{\dimexpr 0.374\linewidth-2\tabcolsep\relax}>{\raggedright\arraybackslash}p{\dimexpr 0.626\linewidth-2\tabcolsep\relax}@{}}
\toprule\noalign{}
Câu hỏi & Trả lời \\
\midrule\noalign{}
\endhead
\bottomrule\noalign{}
\endlastfoot
RQ1: đạt bốn kỳ vọng trên cụm thật? & Có. Hệ thống cuối đạt recall 99,4\%, precision 99,4\%, control 0,6\% và p95 1,97 s trên một campaign mới. Trên LID-DS 2021, cùng tham số, hệ thống phát hiện 77,3\% bản ghi tấn công với 3,0\% bản ghi normal bị báo nhầm; yếu ở các exploit chạy trong chính process của dịch vụ \\
RQ2: feature nào mang tín hiệu? & Cặp exec, kênh syscall tường minh, 5 tổng hợp và cả hai nhánh p-value. Thêm syscall không giúp, và còn làm chậm \\
RQ3: giảm được gánh cảnh báo? & Một phần. Gánh giảm từ 6,74 xuống khoảng 1,5 đợt/replica-giờ nhờ chiết khấu theo replica, nhưng vẫn gấp khoảng 15 lần ngân sách \\
RQ4: tách được process vô hại mới? & ML thì không: 80,7--95,1\% binary vô hại bị gắn cờ. Tripwire tách được lớp tấn công buộc dùng syscall lạ, và gần như không thêm gánh \\
RQ5: RCA chỉ ra process gây ra? & Process tấn công luôn nằm trong 3 nghi phạm đầu. Top-1 14,1\% là do luật phá hoà theo tên; phá hoà theo thời điểm khởi chạy cho 100\% trên cùng lượt (development, chờ kiểm blind) \\
RQ6: chi phí? & Khoảng 250--450 ns mỗi syscall của container được theo dõi, và 13--31\% một core mỗi worker cho phần chấm điểm. Trên ứng dụng thật, probe làm p50 của request tăng 0,5--3\% và p95 tăng 3--8\% ở trung vị; CI quá rộng để nói mức tăng dưới 5\% \\
\end{longtable}

\section{Đóng góp}

\begin{enumerate}
\def\labelenumi{\arabic{enumi}.}
\tightlist
\item
  Một hệ thống chạy live trên Kubernetes, từ kernel tới Alertmanager và cây bằng chứng RCA, với đường chấm điểm khớp từng bit.
\item
  Một cách đánh giá live có đăng ký trước, đối chứng binary vô hại, đo độ trễ và gánh cảnh báo theo giờ.
\item
  Bằng chứng thực nghiệm rằng detector tần suất, kể cả autoencoder, phát hiện sự mới lạ chứ không phải ác ý.
\item
  Chiết khấu bằng chứng theo replica và tripwire A1, cả hai đều đã được xác nhận live.
\end{enumerate}

\section{Hướng phát triển}

\begin{itemize}
\tightlist
\item
  \textbf{Thông tin ngoài tần suất:} tham số syscall và cây process từ Tetragon, để tách hành vi của process mới khỏi việc nó là process mới.
\item
  \textbf{Exploit thật trên cụm:} port một số kịch bản của LID-DS vào một namespace thử nghiệm, gồm tấn công không sinh process mới và tấn công sinh process con từ dịch vụ.
\item
  \textbf{Calibration tự làm mới trong worker:} code đã có. Cần một lượt xác nhận riêng, và một cách loại cửa sổ tấn công không dựa vào lịch campaign.
\item
  \textbf{Calibration theo replica} cho các workload có replica lệch hệ thống, như Kafka.
\item
  \textbf{Thí nghiệm kẻ tấn công thích nghi:} payload chỉ dùng syscall có sẵn, tốc độ dưới 1 lần/s.
\item
  \textbf{RCA:}

  \begin{itemize}
  \tightlist
  \item
    đo độ trễ của nhật ký exec so với case theo sau;
  \item
    luật xếp hạng tách được hai process mới chạy gần nhau;
  \item
    xếp hạng có hỗ trợ của LLM trên bằng chứng thật, theo hướng của CORTEX (\citealp{wei2025}).
  \end{itemize}
\item
  \textbf{Đưa một phần việc chấm điểm vào kernel} bằng số cố định (\citealp{osaki2024}), để giảm độ trễ và chi phí user space.
\end{itemize}
